\chapter{Aspectos de Implementação e Avaliação}
\label{ch:implementacao}

%%%%%%%%%%%%%%%%%%%%%%%%%%%%%%%%%%%%%%%%%%%%%%%%%%%%%%%%%%%%

\section{Tecnologias Utilizadas}
\label{sec:tecnologias}

O protótipo do eProfile será desenvolvido em Java \cite{java}. O Java é uma linguagem de código aberto, orientada a objetos e de alto nível, desenvolvida com o intuito de poder ter seus programas executados em qualquer arquitetura.

As bases de dados serão implementadas utilizando o banco de dados SQLite3 \cite{sqlite3}. O SQLite3 é um banco de dados SQL auto-contido transacional, disponível em domínio público. O SQLite3 foi escolhido por apresentar bom desempenho e não necessitar de um servidor separado para sua implantação.

%%%%%%%%%%%%%%%%%%%%%%%%%%%%%%%%%%%%%%%%%%%%%%%%%%%%%%%%%%%%

\subsection{Motor de Inferência}
\label{sec:motor_inferencia}

Foram estudados os seguintes raciocinadores que oferecem suporte a estas operações: Pellet \cite{sirin05}, FaCT++ \cite{tsarkov06} e HermiT \cite{motik09}.

\begin{itemize}

\item {\bf Pellet:} é um racionador com suporte a OWL DL, que oferece bom desempenho e um \emph{middleware} extensível. O Pellet é escrito em Java pode ser utilizado para raciocinar com tipos de dados definidos pelo usuário e provê suporte para depuração em ontologias \cite{sirin05}.

\item {\bf FaCT++:} é um raciocinador de lógica proposicional criado como uma plataforma para experimentar com novos algoritmos de \emph{tableaux} semântico e técnicas de otimização. Ele incorpora a maior parte das técnicas de otimização padrão, mas inclui algumas novas \cite{tsarkov06}.

\item {\bf HermiT:} é o primeiro raciocinador OWL que se baseia em um cálculo \emph{hypertableaux}, que provê desempenho superior a qualquer algoritmo anterior \cite{motik09}. O HermiT é disponibilizado como um plugin do Protegé, como uma ferramenta de linha de comando ou como uma biblioteca Java.

\end{itemize}

Devido ao conjunto de funcionalidades e ao desempenho superior, o HermiT foi escolhido como o raciocinador utilizado pelo eProfile.

%%%%%%%%%%%%%%%%%%%%%%%%%%%%%%%%%%%%%%%%%%%%%%%%%%%%%%%%%%%%

\subsection{Linguagem de Ontologia}
\label{sec:ling_ontologia}

Para o armazenamento e processamento das ontologias, foram estudadas as linguagens de representação de conhecimento CycL \cite{lenat89}, RDF/S \cite{rdfs}, DAML+OIL \cite{daconta03} e OWL \cite{daconta03}.

% CycL
A linguagem \textbf{CycL} é a linguagem utilizada pelo o projeto de \emph{Cyc} \cite{lenat89}. O Cyc é um projeto de inteligência artificial que tem o objetivo de montar uma ontologia e base de conhecimento a partir do senso comum, com o propósito de permitir que aplicações de inteligência artificial realizem raciocínio semelhante ao humano. A linguagem CycL é uma linguagem declarativa baseada em lógica de primeira ordem clássica, com extensões para operadores modais e lógica de ordem superior. A base de conhecimento é dividida em micro-teorias (Mt), que são coleções de conceitos e fatos tipicamente pertencentes a um domínio particular de conhecimentos e, ao contrário da base de conhecimento como um todo, cada micro-teoria deve estar livre de contradições.

% RDF/S
\textbf{RDF} (\emph{Resource Description Framework}) é uma linguagem utilizada para representação de metadados a respeito de recursos na Web \cite{rdf}. No entanto, o conceito de ``recurso da Web'' pode ser generalizado, de forma que o RDF pode ser utilizado para representar informações a respeito de coisas que podem ser identificadas na Web, mesmo que elas não possam ser carregadas da Web. Recursos podem ser recursos eletrônicos, como arquivos, ou conceitos, como pessoas. O RDF é complementado pelo \textbf{RDF/S} (\emph{RDF/Schema}), uma extensão semântica do RDF. O RDF/S é uma linguagem de descrição de vocabulário que provê mecanismos para descrever grupos de recursos relacionados e os relacionamentos entre estes recursos \cite{rdfs}.

% DAML+OIL
O \textbf{DAML+OIL} é uma linguagem de representação de conhecimento baseada nas linguagens DAML e OIL. O DAML (\emph{DARPA Agent Markup Language}) foi desenvolvido pelo Centro de Pesquisas Avançadas da Defesa Americana e é uma linguagem de marcação de agentes baseada no RDF. Já o OIL (\emph{Ontology Inference Layer}) é uma linguagem com os mesmos objetivos, criada pelo por um grupo de universidades europeias. Ambos os esforços foram unidos na criação do DAML+OIL, que por sua vez deu origem ao OWL \cite{daconta03}.

% OWL
O \textbf{OWL} é a linguagem de ontologia mais expressiva definida ou em definição para a Web Semântica \cite{daconta03}. O OWL utiliza uma URI (\emph{Uniform Resource Identifier}) para os nomes, e a estrutura de descrição disponibilizada pelo RDF. O OWL está uma camada acima do RDF e do RDF/S, e adiciona mais vocabulário para descrever propriedades e classes: entre outros, relações entre classes, cardinalidade, igualdade, tipagem mais avançada de propriedades, características de propriedades e classes enumeradas \cite{owl-guide}. O OWL está disponível em três sublinguages: \emph{OWL Lite}, \emph{OWL DL} e \emph{OWL Full}, cada uma aumentando o nível de complexidade e expressividade. Um grande número de ferramentas com capacidade para edição e raciocínio semântico com suporte para OWL está disponível no mercado.

% Escolhido - OWL
Devido à expressividade e ao suporte de ferramentas externas, o OWL foi escolhido como a linguagem de ontologias utilizado pelo eProfile para modelagem da ontologia do perfil e das trilhas.

%%%%%%%%%%%%%%%%%%%%%%%%%%%%%%%%%%%%%%%%%%%%%%%%%%%%%%%%%%%%

\section{Avaliação}
\label{sec:avaliacao}

\todo[inline]{Detalhar melhor esta seção}

A validação deste modelo busca provar que o modelo apresentado é funcional e efetivo.

Para o teste de funcionalidade, o modelo será implementado e integrado com o eTrail, um gerenciador de trilhas. Uma vez que esta integração esteja realizada, um conjunto de regras será desenvolvido, de modo a verificar a capacidade de inferência do eProfile.

A seguir, será testada a capacidade do eProfile em gerenciar perfis de forma efetiva. Para tal, um grupo de usuários utilizará um aplicativo cujo perfil é gerenciado pelo eProfile. Após a utilização, será apresentado a estes usuários um questionário que buscará avaliar a efetividade do eProfile.
